\section{System Design}
\label{sec:design}

\system turns the three observations into three decisions: which Agents
receive a recoverable checkpoint (O1), which regions of their KV state
receive fast replicas (O2), and how much HBM those replicas may occupy
(O3). The objective is to reduce workflow completion time after a worker
failure while bounding the memory and transfer resources consumed by
protection during normal execution.

\subsection{Overview and Recovery State}
\label{sec:teacher-overview}

The design operates on DAGs of LLM invocations and their accepted token
histories. It uses two managed memory tiers. Fault-safe DRAM retains a
complete checkpoint for each admitted Agent; a compatible takeover worker
keeps selected replicas of that checkpoint in Peer HBM. DRAM provides
recovery coverage, while Peer HBM removes selected restoration stalls.
An Agent without an admitted checkpoint can still recover by prefilling
its retained input and accepted output tokens.

Figure~\ref{fig:teacher-overview} shows how the components cooperate.
The workflow runtime supplies invocation dependencies and progress.
The \emph{Workflow Tracker} estimates remaining work and slack.
The \emph{Recovery-aware Router} dispatches requests using load and prefix
locality and identifies compatible takeover workers outside the source
failure domain. The \emph{KV Master} selects DRAM checkpoints and Peer
regions using the current workflow and recovery profiles. A per-worker
\emph{Adapter} publishes a safe Peer budget; rank-local \emph{Connectors}
execute Master's copies and report their completion. KV bytes move
between memory resources, while these components exchange metadata.

\begin{figure}[t]
 \centering
 \includegraphics[width=\columnwidth]{figures/teacher-system-overview.pdf}
 \caption{Two-tier protection organized around O1--O3. Each worker is a
 DP replica with two TP ranks. DRAM checkpoints contain all layers and
 shards; Peer regions are additional copies on corresponding ranks of a
 compatible takeover worker. DRAM and accepted histories survive the
 protected source failure. Blue paths carry KV; dashed paths carry
 progress, plans, completion status, or capacity feedback.}
 \label{fig:teacher-overview}
 \Description{Application, Controller, and Engine bands. Tracker feeds
 workflow estimates to KV Master; Router selects execution and takeover
 workers. Master exchanges plans and status with Connectors, while
 Adapter reports a Peer quota. Two workers have distinct TP0 and TP1
 execution and Peer pools, with rank-matched Peer copy paths and a
 complete fault-safe DRAM checkpoint service outside their boundaries.}
\end{figure}

For an Agent $a$, let $T_a$ be the length of its accepted token history
and $M_a\le T_a$ the latest complete checkpoint frontier. Recovery restores
the prefix through $M_a$ and replays the accepted tail
$N_a=T_a-M_a$. A checkpoint contains every required layer, TP shard,
and attention auxiliary object for that prefix. Model revision, numerical
format, positional metadata, KV layout, and prefix identity must match
before any object can be reused. This extends checkpoint-and-replay
recovery~\cite{strati2024dejavu,cao2026lumen} with selective fast residency.

Let $D_a$ denote the DRAM objects for this complete checkpoint, and let
$H_a$ be the logical regions additionally copied into Peer HBM. The
invariant is $H_a\subseteq D_a$: a Peer copy never replaces its DRAM
backing. Execution KV belongs to the serving engine and is accounted
separately. Physical DRAM objects shared by compatible prefixes are
charged once, but every additional Peer, staging, or in-flight copy
consumes capacity. The advertised DRAM service, token history, and
takeover worker must remain reachable after the protected source failure;
DRAM on the same host does not provide host-failure protection.

\subsection{O1: Workflow-aware Agent Protection}
\label{sec:teacher-agent-protection}

O1 requires the system to value a checkpoint by the workflow delay it
avoids. Tracker first constructs the visible residual invocation DAG.
Completed nodes contribute no remaining work, and already running nodes
contribute only their unfinished execution. An Agent may issue several
invocations; dependency edges connect those calls to the rest of the
workflow. The current snapshot does not invent unseen future branches.

\paragraph{Estimating urgency.}
A GBDT estimates each invocation's remaining duration from information
available at the snapshot: model and Agent role, current input length,
decoded-token progress, elapsed execution, and lengths and timings of
completed calls. Training and calibration separate whole workflows;
future outputs are never input features. Forward and backward
critical-path passes~\cite{kelley1959cpm} compute earliest and latest
finishes, $EF_v$ and $LF_v$. Tracker publishes guarded slack
\begin{equation}
 s_v=\max\{0,LF_v-EF_v-g_v\},
 \label{eq:teacher-slack}
\end{equation}
where $g_v$ is a calibrated upper quantile of positive slack
overestimation. Requests on the same residual path consume a shared
margin; slack is not a separate allowance that can be spent repeatedly
at every invocation. These estimates are advisory. Invalid or stale
inputs grant no positive slack, and unsupported DAGs use the fallback
priority described below. Slack guides urgency and queue tie-breaking;
the joint workflow calculation determines the final replacement.

\paragraph{Valuing a common failure.}
For a placement $P$, a failure of source worker $f$ interrupts the entire
set of invocations assigned to that worker. Router fixes compatible
takeover targets for a planning snapshot. A recovery profiler schedules
reconstruction or restoration and accepted-tail replay as prerequisite
tasks of each interrupted invocation; only its remaining normal work
follows them. If a later visible invocation needs history lost with the
worker, its restoration is likewise scheduled before that invocation.
The schedule includes concurrent recoveries, target
admission, and actual shared transfer links. Streams
contend only when they share a resource; the model does not impose a
single global recovery queue. The same ready-time queue rule, with
guarded slack and request ID breaking ties, is used for every candidate.

Let $F_w^f(P)$ be the predicted completion time of workflow $w$ under
this resource schedule and its residual DAG, and $F_w^0$ its finish
without the failure in the same snapshot. Surviving requests delayed by
shared recovery resources are included. Master compares
\begin{equation}
 J(P)=\sum_f \pi_f\sum_w\bigl(F_w^f(P)-F_w^0\bigr),
 \label{eq:teacher-objective}
\end{equation}
where $\pi_f$ are fixed, normalized source-failure weights for the
comparison; equal weights are the default. They are not measured fault
probabilities. Each candidate uses the same workload snapshot, queue
policy, and failure weights. At a join, finish time is determined by the
maximum predecessor finish, so this calculation preserves O1's
non-additive effect instead of summing isolated slack penalties.

\paragraph{Admitting complete checkpoints.}
DRAM admission chooses among complete sealed frontiers for each Agent,
including the no-checkpoint option $M_a=0$. A candidate can admit a new
checkpoint, advance an existing frontier, or retire a checkpoint whose
capacity is more valuable elsewhere. A retiring checkpoint first loses
its Peer replicas; its token history remains available for recomputation.
For shared physical objects, a victim releases bytes only after the last
retained manifest stops referencing them.

Master evaluates single-Agent changes and small joint changes. In
particular, it follows the currently limiting paths into each exposed
join and offers one bundle containing the available checkpoint upgrades
for all affected co-bottleneck Agents. For example, if either of two
branches remains 10\,s late unless both are protected, their two-Agent
bundle can advance the join even when each singleton has zero value.
This is a candidate-generation rule, not an assumption that benefits
become additive. Bounded search can still miss other combinations.

When DRAM is full, a candidate includes its victims. Master tentatively
removes the eligible checkpoint with the smallest recomputed increase
in $J$ per released physical byte, repeating until the entire candidate
fits. It then compares the whole replacement with the current plan.
An optional change requires $J(P)-J(P')>\epsilon$, a configured
hysteresis margin, and feasible copy and memory reservations. Ties keep
the existing placement. Per-round work and transfer-byte limits bound
decision overhead and churn. If no candidate fits, the previous
checkpoint remains, or the Agent uses recomputation. The admission
decision never advertises a partially copied checkpoint as complete.

If Tracker cannot supply a valid residual DAG or timing profile, Master
uses profiled request-recovery time saved per byte and deterministic
ties. It retains the same capacity and commit rules, but makes no
workflow-priority claim for this fallback. Invocation-level labels and
held-out calibration are required before positive slack can guide
protection.

\subsection{O2: Selecting KV Regions by Exposed Recovery Delay}
\label{sec:teacher-region-selection}

For an Agent admitted to DRAM, O2 asks which parts of its checkpoint are
worth an additional Peer copy. A managed region identifies a layer and
a historical token block together with its required TP objects. Peer
placement maps each shard to the corresponding rank of the takeover
worker; budgets are checked per rank. Auxiliary state needed to execute
the native selector or attention kernel is tracked as a dependency,
rather than assumed to be free.

\paragraph{Causal demand candidates.}
The engine reports native attention selections over a recent completed
window. For region $b$ in layer $\ell$, Master estimates
\begin{equation}
 \widehat p_{a,\ell,b}=
 \frac{\text{eligible selection units that selected }b}
      {\text{selection units for which }b\text{ was visible}}.
 \label{eq:teacher-demand}
\end{equation}
The selection unit follows the model's native rule, rather than a
kernel tile chosen for execution. Head duplicates count once within a
unit; causal visibility is determined before sparse pruning. Small
denominators use a calibrated layer prior. For dense attention,
comparable visible historical blocks have access factor one.
Historical frequency supplies a candidate ordering, not a guarantee
about replay demand or the benefit of a Peer copy.

Master considers individual regions and progressively larger hot sets
within a layer. Mandatory state accompanies every candidate that needs
it. Recent consecutive selection windows provide demand scenarios for
profiling; keeping them together preserves correlated accesses and
native reveal order. Both the window and the profiling budget are fixed
configuration parameters. Missing or invalid sparse statistics use
full-layer transfers as a conservative fallback. Protection never
changes the model's native attention selections.

\paragraph{Accounting for overlap.}
The recovery profile executes the accepted tail against the saved
prefix. It models transport from DRAM, bounded prefetch buffers, native
demand misses, and layer computation on the same timeline. Prefetches
can start early; a miss cannot be requested before its native selection
reveals the demand. Let $q_\ell$ be the time layer $\ell$ is ready to
consume historical KV, including query and selector preparation, and
let $a_\ell(H_a)$ be the last required arrival under Peer set $H_a$.
With remaining layer compute time $c_\ell$, the profiler uses
\begin{equation}
 \begin{aligned}
 e_\ell(H_a)&=[a_\ell(H_a)-q_\ell]^+,\\
 f_\ell(H_a)&=q_\ell+e_\ell(H_a)+c_\ell.
 \end{aligned}
 \label{eq:teacher-overlap}
\end{equation}
Later readiness depends on earlier finishes. Transfer completion times
also depend on link sharing and the finite prefetch schedule, so a
candidate must rerun the affected timeline. The opportunity to hide
transfer comes from remaining replay computation, not computation
already finished before the failure.

For each candidate hot set, this profile produces a new recovery delay,
which is inserted into Equation~\ref{eq:teacher-objective}. A frequently
used block can have zero marginal benefit when its read is already
hidden. Conversely, removing an exposed early-layer read can advance
several later layers. Master chooses the feasible candidate with the
largest positive reduction in $J$ per additional Peer byte, recomputes
after each tentative change, and stops at the budget or work limit.
Layer bundles and the co-bottleneck bundles from O1 allow jointly useful
changes to be considered. This bounded heuristic claims neither
globally optimal placement nor a universally optimal layer depth.

Peer replacement evaluates its victims in the same comparison. If one
rank cannot accommodate the required shards, the candidate waits rather
than counting an incomplete cross-rank acceleration. A common snapshot
fixes sampled demand windows across candidate evaluations. Prediction
error can increase waiting, but every demand miss remains recoverable
from the complete DRAM checkpoint.

\begin{figure}[t]
 \centering
 \includegraphics[width=\columnwidth]{figures/teacher-selective-recovery.pdf}
 \caption{Complete DRAM coverage and selective Peer acceleration.
 The illustrative checkpoint has three layers and four historical
 blocks; every cell represents its required TP shards. Peer holds
 additional copies of L1/B0, L1/B2, and L2/B1. When L1 needs B0, B1,
 and B2, the two Peer hits combine with B1 fetched from DRAM.
 The accepted tail is replayed. Peer eviction discards only a replica
 after readers drain; it neither removes DRAM coverage nor requires
 writing the checkpoint back. This is a state diagram, not a timing result.}
 \label{fig:teacher-selective-recovery}
 \Description{A complete three-by-four DRAM grid and a separate Peer
 grid containing three highlighted replicas. An uncommitted tail is
 outside both checkpoint grids. Recovery consumes resident B0/B2 and a
 DRAM fetch of B1, then replays accepted tail tokens.}
\end{figure}

\subsection{O3: Adapting the Peer HBM Budget}
\label{sec:teacher-capacity}

O3 separates content selection from capacity authorization. Master
chooses which regions occupy Peer; Adapter determines how much Peer
memory foreground execution can safely support. For takeover rank $k$,
let $B^{\rm pool}_k$ exclude weights and fixed allocations,
$B^{\rm serve}_k$ cover the serving high watermark and admitted execution
reservations, and $B^{\rm reserve}_k$ reserve growth and workspace.
The physical ceiling is
\begin{equation}
 C_k=\max\{0,\min(B^{\rm cap}_k,
 B^{\rm pool}_k-B^{\rm serve}_k-B^{\rm reserve}_k)\}.
 \label{eq:teacher-ceiling}
\end{equation}
Peer occupancy is not included in $B^{\rm serve}_k$ and is not subtracted
twice. Every committed, reserved, copying, and retiring Peer byte counts
against physical use. Distinct TP ranks have distinct ceilings.

Adapter uses additive increase and multiplicative
decrease~\cite{chiu1989aimd}. On allocation stalls, reserve violations,
or memory-related admission waits, it multiplies the current quota
$B_k$ by $0<\gamma<1$ and clips it to $C_k$. After $m$ consecutive
stable windows below a lower wait threshold, it adds $\Delta$ bytes
only if Master reports a beneficial candidate deferred for capacity.
This demand signal is recorded before quota filtering, so a zero quota
does not prevent later growth. Every update is page-aligned and clipped
to the current ceiling; missing telemetry disables growth. The gap
between pressure and stable thresholds suppresses rapid oscillation.

When the budget shrinks, Master stops new Peer admissions and retires
the eligible replicas with the smallest recomputed workflow penalty
per byte until occupancy fits the new target. Mandatory shrink does
not require a positive improvement in $J$. The complete DRAM backing
makes this a discard operation, with no additional checkpoint offload.
Reader and copy leases still delay physical release. Foreground
admission uses actual occupancy and reservations; it cannot spend
memory merely because a lower quota was published. Under immediate
pressure, pending Peer copies can be cancelled and new foreground
allocations wait until enough physical space is available.

\subsection{Online Checkpointing and Recovery}
\label{sec:teacher-execution}

The decisions above run at request arrival and completion, DAG/progress
updates, memory-pressure events, and a fixed block-aligned checkpoint
interval $Z$. Prefill offers sealed prefix blocks; decode offers new
sealed state every $Z$ accepted tokens and at final request completion.
The interval controls review frequency, not the guaranteed replay tail:
copy lag can leave $N_a>Z$. Only completed native observations update
the next planning snapshot.

Algorithm~\ref{alg:teacher-protection} summarizes the common event loop.
\textsc{Candidates} produces the single-Agent, layer, and co-bottleneck
changes described above. \textsc{FitWithVictims} adds any required
retirements and recomputes their joint cost; its reservations include
both old and new versions during copying. Benefit-per-byte ranking is
used within one memory pool, not to treat a DRAM byte and a Peer byte
as interchangeable. A bounded pass over DRAM precedes Peer selection;
later events reconsider the resulting placement. \textsc{ReserveIfCurrent}
checks graph, manifest, and quota versions as well as transient memory
and copy limits. A stale or infeasible plan is requeued without publishing
its proposed state.

\par\medskip\noindent\begin{minipage}{\columnwidth}
\raggedright
\refstepcounter{designalgorithm}\label{alg:teacher-protection}
\textbf{Algorithm \thedesignalgorithm: Observation-driven protection}
\par\smallskip
\begin{tabularx}{\columnwidth}{@{}r@{\quad}X@{}}
\toprule
1 & \textsc{ReconcileCopies}$(e)$; snapshot progress and budgets \\
2 & $P\gets$ current committed placement \\
3 & Retire expired state; apply mandatory Peer shrink \\
4 & \textbf{for} pool $j\in[\mathrm{DRAM},\mathrm{Peer}]$ \textbf{do} \\
5 & \quad $Q\gets\textsc{Candidates}(j,P)$ \\
6 & \quad \textbf{while} work budget remains \textbf{do} \\
7 & \qquad For each examined $c\in Q$, form \\
8 & \qquad $P_c\gets\textsc{FitWithVictims}(P,c,j)$ \\
9 & \qquad Recompute $J(P)-J(P_c)$ on one snapshot \\
10 & \qquad $P^*\gets$ best feasible gain/byte candidate \\
11 & \qquad \textbf{if} none exceeds $\epsilon$ \textbf{then break} \\
12 & \qquad $P\gets P^*$; refresh affected candidates \\
13 & \textbf{if not} \textsc{ReserveIfCurrent}$(P)$ \textbf{then} \\
14 & \quad \textsc{Requeue}$(e)$; \textbf{return} \\
15 & Submit DRAM commits, then dependent Peer copies \\
16 & Copy completion or cancellation emits a new event \\
\bottomrule
\end{tabularx}
\end{minipage}\par\medskip

\paragraph{Publishing a checkpoint.}
Connector copies immutable sealed objects to the advertised DRAM
domain. Master advances $M_a$ only after every required layer, TP shard,
and auxiliary object is confirmed, then atomically publishes the new
manifest. Old objects remain charged until their references and readers
drain. Peer installation starts from a committed manifest, and the
directory publishes a hit only after the corresponding target copy
finishes. Source failure during an incomplete update leaves the prior
committed frontier usable. If temporary space is insufficient, the
update is deferred rather than releasing the old checkpoint early.

\paragraph{Resuming an interrupted invocation.}
On confirmed failure, Router fences the old execution epoch and
activates a compatible surviving target. Connector acquires a lease on
the latest complete DRAM manifest, consumes valid Peer hits, and fetches
missing regions under the bounded transfer schedule. It replays the
accepted tail using the exact recorded token sequence and resumes
generation without emitting accepted tokens twice. If the preferred
target is unavailable, another compatible target restores from DRAM;
Peer residency on the old target is not counted as a local hit.

Checkpoint retirement stops new readers and invalidates dependent Peer
replicas before releasing DRAM. Existing recovery leases keep their
objects alive. If no valid checkpoint remains, the surviving input and
token log define the explicit full-prefill fallback. Thus priority and
demand estimates affect recovery time, while complete manifests,
version matching, and reader lifetimes determine which state can safely
be consumed. The mechanisms are a design specification; integrated
fault-injection measurements are still required to establish WCT gains,
prediction quality, and normal-service overhead.
